\begin{examplebox}
	\textbf{\textcolor{CExample}{例1（基础）}}
	\textbf{\textcolor{CExample}{题目：}}判断直线$l_1:2x+3y-5=0$与$l_2:3x-2y+1=0$是否垂直。\textbf{\textcolor{CExample}{解题步骤：}}
	\begin{description}[style=nextline, leftmargin=2.8em, labelsep=0.5em, itemsep=0.45em]
		\item[\textbf{第一步（识别系数）：}]
			从方程中读取系数：$A_1=2,B_1=3$；$A_2=3,B_2=-2$。
			\[
				A_1=2,\;B_1=3;\quad A_2=3,\;B_2=-2.
			\]
			\par\textbf{理由：}直线的一般式方程中直接读取。
		\item[\textbf{第二步（计算点积）：}]
			计算$A_1A_2+B_1B_2$：
			\[
				\begin{aligned}
					A_1A_2+B_1B_2 &= 2\times3+3\times(-2) \\ &= 6-6 \\ &= 0.
			\end{aligned}
			\]
			\par\textbf{理由：}算术运算。
		\item[\textbf{第三步（判定垂直）：}]
			由于$A_1A_2+B_1B_2=0$，根据垂直充要条件得$l_1\perp l_2$。
			\[
				l_1\perp l_2.
			\]
			\par\textbf{理由：}两直线垂直的充要条件是$A_1A_2+B_1B_2=0$。
	\end{description}
	\textbf{\textcolor{CExample}{关键结论：}}
	\textbf{因为$A_1A_2+B_1B_2=0$，所以$l_1\perp l_2$。}

	\medskip
	\textbf{\textcolor{CExample}{例2（稍变形）}}
	\textbf{\textcolor{CExample}{题目：}}已知直线$l_1:(m+1)x+2y-4=0$与$l_2:x+my+1=0$垂直，求实数$m$的值。\textbf{\textcolor{CExample}{解题步骤：}}
	\begin{description}[style=nextline, leftmargin=2.8em, labelsep=0.5em, itemsep=0.45em]
		\item[\textbf{第一步（写出系数）：}$A_1=m+1,B_1=2$；$A_2=1,B_2=m$。
			\[
				A_1=m+1,\;B_1=2;\quad A_2=1,\;B_2=m.
			\]
			\par\textbf{理由：}对照一般式直接得到。
		\item[\textbf{第二步（代入垂直条件）：}]
			由$A_1A_2+B_1B_2=0$得$(m+1)\times1+2\times m=0$，即$m+1+2m=0$。
			\[
				(m+1)+2m=0.
			\]
			\par\textbf{理由：}垂直充要条件。
		\item[\textbf{第三步（解方程）：}]
			整理得$3m+1=0$，解得$m=-\frac13$。
			\[
				\begin{aligned}
					3m+1 &= 0 \\ m &= -\frac13.
			\end{aligned}
			\]
			\par\textbf{理由：}解一元一次方程。
	\end{description}
	\textbf{\textcolor{CExample}{关键结论：}}
	\textbf{$m=-\dfrac13$。}

	\medskip
	\textbf{\textcolor{CExample}{例3（综合应用）}}
	\textbf{\textcolor{CExample}{题目：}}求过点$P(1,2)$且与直线$3x-4y+7=0$垂直的直线方程。\textbf{\textcolor{CExample}{解题步骤：}}
	\begin{description}[style=nextline, leftmargin=2.8em, labelsep=0.5em, itemsep=0.45em]
		\item[\textbf{第一步（设直线系数）：}]
			设所求直线为$A_2x+B_2y+C_2=0$。由垂直条件得$3A_2+(-4)B_2=0$。
			\[
				3A_2-4B_2=0.
			\]
			\par\textbf{理由：}垂直充要条件$A_1A_2+B_1B_2=0$。
		\item[\textbf{第二步（取一组系数）：}]
			取$A_2=4,B_2=3$，满足$3\times4-4\times3=0$。设直线为$4x+3y+C_2=0$。
			\[
				4x+3y+C_2=0.
			\]
			\par\textbf{理由：}满足垂直条件的系数不唯一，通常取最简整数比。
		\item[\textbf{第三步（求常数项）：}]
			将点$P(1,2)$代入$4x+3y+C_2=0$，得$4\times1+3\times2+C_2=0$，即$4+6+C_2=0$，解得$C_2=-10$。因此所求直线为$4x+3y-10=0$。
			\[
				4x+3y-10=0.
			\]
			\par\textbf{理由：}直线过点，则点坐标满足方程。
	\end{description}
	\textbf{\textcolor{CExample}{关键结论：}}
	\textbf{所求直线方程为$4x+3y-10=0$。}
\end{examplebox}
